\documentclass{article}
\usepackage[utf8]{inputenc}
\usepackage{polski}
\usepackage{geometry}
\geometry{a4paper, margin=1in}

\title{Raport Techniczny: Prototyp Systemu Inspekcji Wizualnej}
\author{Zespół Projektowy}
\date{\today}

\begin{document}

\maketitle

\section{Założenia Projektowe}
Celem projektu jest stworzenie prototypu systemu do wizualnej inspekcji jakości. Fundamentem systemu jest model głębokiego uczenia, który potrafi odróżnić produkty wadliwe od poprawnych. W ramach obecnego etapu, skupiamy się na dokończeniu implementacji autoenkodera konwolucyjnego (CAE) i zintegrowaniu go z prostą aplikacją webową, która posłuży jako demonstracja jego możliwości.

\section{Model Detekcji Anomalii: Autoenkoder Konwolucyjny (CAE)}
Sercem naszego systemu detekcji jest autoenkoder konwolucyjny. Model ten jest trenowany wyłącznie na obrazach produktów bez wad. Jego zadaniem jest nauczenie się efektywnej, skompresowanej reprezentacji danych. Podczas inspekcji, obrazy o wysokim błędzie rekonstrukcji (tj. te, których model nie potrafi poprawnie odtworzyć) są klasyfikowane jako anomalie.

\section{Aplikacja Webowa i Backend}
Aby umożliwić interakcję z naszym modelem, tworzymy aplikację webową składającą się z backendu i prostego interfejsu użytkownika.
\begin{itemize}
    \item \textbf{Backend API:} Zdecydowaliśmy się na stworzenie backendu w oparciu o FastAPI. Udostępni on kluczowy endpoint \texttt{/infer}, który będzie przyjmował obraz w formacie base64, przekazywał go do modelu, a następnie zwracał wynik w postaci wskaźnika anomalii.
    \item \textbf{Interfejs użytkownika:} Do szybkiej wizualizacji działania modelu wykorzystamy Gradio. Umożliwi to łatwe przesyłanie obrazów i natychmiastowe otrzymywanie wizualnej informacji zwrotnej bez potrzeby budowania złożonego frontendu.
\end{itemize}

\section{Kluczowe Technologie}
Wybór technologii został podyktowany szybkością rozwoju i dostępnością narzędzi w ekosystemie.
\begin{itemize}
    \item \textbf{Python:} Wszechstronny język z bogatym ekosystemem do uczenia maszynowego i tworzenia aplikacji webowych.
    \item \textbf{PyTorch:} Elastyczny framework do głębokiego uczenia, idealny do prototypowania i budowy niestandardowych modeli, takich jak nasz autoenkoder.
    \item \textbf{FastAPI:} Nowoczesny i wydajny framework webowy dla Pythona, który umożliwia szybkie tworzenie API, co jest kluczowe dla naszego backendu.
    \item \textbf{Gradio:} Biblioteka do szybkiego tworzenia interfejsów użytkownika dla modeli uczenia maszynowego, co pozwala nam na szybką demonstrację działania naszego autoenkodera.
    \item \textbf{Docker:} Narzędzie do konteneryzacji, które zapewni spójne i odizolowane środowisko uruchomieniowe dla naszej aplikacji, co ułatwi jej wdrożenie.
\end{itemize}

\section{Kierunki Dalszego Rozwoju}
Po pomyślnym wdrożeniu obecnej wersji, planujemy rozszerzyć system o następujące elementy:
\begin{itemize}
    \item Zintegrowanie alternatywnych, bardziej zaawansowanych modeli detekcji anomalii, takich jak PaDiM, w celu porównania ich wydajności.
    \item Rozbudowa interfejsu użytkownika o możliwość przeglądania historii inspekcji oraz bardziej szczegółowe wizualizacje.
    \item Dodanie bazy danych do przechowywania wyników inspekcji i metadanych obrazów.
    \item Stworzenie modułu do generowania raportów z przeprowadzonych kontroli jakości.
\end{itemize}

\end{document}
